\section{BH2: The cost of trading in the window}\label{sec:bh2}\label{sec:book}

The second hypothesis is that the window is more costly to trade in on expiry days. Its primary
measure is the quoted spread. 
On the \DesignSessions{} sessions of the order book sample the limit order book is
reconstructed, which shows how the supply of liquidity in the window changes on expiry days.
Each measure is computed per security and session and compared between each monthly expiry
and the control session in the same month, the trading day before it.
Table~\ref{tab:descriptives} reports the averages and Table~\ref{tab:contrasts} compares the
expiry-minus-control differences of the securities with derivatives with those of the
placebo group. The full set of paired tests is reported in Appendix~\ref{app:tests}.

\input{generated/tables/descriptives}

The window is more expensive to trade in on expiry days. The quoted spread widens from
\SpreadMeanBpsControl{} to \SpreadMeanBpsExpiry{} basis points, the displayed notional in the
book falls from \BookVisibleCrControl{} to \BookVisibleCrExpiry{} crore, and sweeping the
displayed book with an order of \ProbeSizeCr{} crore moves the price by
\MoveBpsSmallExpiry{} basis points against \MoveBpsSmallControl{} on control sessions. The share
of orders cancelled within one second of entry rises from \LifePhantomRateControl{}\% to
\LifePhantomRateExpiry{}\%. The widening of the spread and the rise in fleeting orders are
specific to securities with derivatives: relative to the placebo group, the spread widens by an
additional \DidMeanspreadbpsPlacebo{} basis points ($p = \DidMeanspreadbpsPlaceboP{}$) and the
share of orders cancelled within a second rises by an additional
\DidPhantomorderratePlaceboPp{} percentage points ($p = \DidPhantomorderratePlaceboP{}$). The
cost of sweeping the book rises by a similar amount in both sets of securities.

The rise in cost is modest relative to the level. Sweeping the book is several times more
expensive in the illiquid group, at \MoveBpsSmallIlliquid{} basis points, than in the liquid
group, at \MoveBpsSmallLiquid{}, on any session. The cost of moving the price is set far more by
the security than by the day.


The primary test of BH2 is rejected: the quoted spread is wider on expiry sessions
($d = \HOneTwoEffectSizeCohenD{}$, $p = \HOneTwoPValue{}$), and the widening is specific to
securities with derivatives. The supporting measures agree. Displayed depth near the touch
falls ($p = \HOneFourPValue{}$), the price move from a given flow over ten seconds rises
($d = \PFourEffectSizeCohenD{}$), realised variance per minute rises relative to the rest of the
session ($p = \HTwoFourPValue{}$) and so does illiquidity measured as price change per unit of
turnover ($p = \HTwoSevenPValue{}$). The price impact of a single trade does not change
($p = \HOneSevenPValue{}$), and the spread recovers from shocks no more slowly
($p = \HThreeZeroPValue{}$).

Figure~\ref{fig:profile} follows the order book through the window minute by minute, and
Figure~\ref{fig:cost} compares the cost of sweeping the book across groups. The expiry and
control sessions differ throughout the window rather than only at its end.

\begin{figure}[htbp]
  \centering
  \includegraphics[width=\textwidth]{window_profile.pdf}
  \caption{Quoted spread, share of window volume, concealed share of resting size and
    cancellations per entry by minute of the window, for expiry and control sessions of each
    group. Order book sample.}
  \label{fig:profile}
\end{figure}

\begin{figure}[htbp]
  \centering
  \includegraphics[width=0.85\textwidth]{marking_cost.pdf}
  \caption{Price move from sweeping the displayed book with \ProbeSizeCr{}~crore, and the
    displayed notional in the book, by group, on expiry and control sessions. Medians over
    snapshots. The sweep assumes the book neither replenishes nor reacts.}
  \label{fig:cost}
\end{figure}

\subsection{The book through the expiry day}

The snapshots of the order book cover the whole session, which shows whether the cost of
trading rises only in the window or across the expiry day. Each security's quoted spread and
displayed depth within five price levels are averaged by minute, and the expiry day is compared
with the day before it, security by security, as a difference in logarithms.

\begin{figure}[htbp]
  \centering
  \includegraphics[width=\textwidth]{bh2_spread_day.pdf}\\[0.5em]
  \includegraphics[width=\textwidth]{bh2_depth_day.pdf}
  \caption{Quoted spread (top) and displayed depth within five price levels (bottom) through
    the session on monthly expiries and on the day before, medians across securities with a
    five-minute moving average, logarithmic scale. The settlement window is shaded.}
  \label{fig:bookday}
\end{figure}

\begin{figure}[htbp]
  \centering
  \includegraphics[width=\textwidth]{bh2_window_zoom.pdf}
  \caption{Quoted spread and displayed depth minute by minute over the last hour of the
    session, for the securities with derivatives and the placebo group. Medians across
    securities.}
  \label{fig:windowzoom}
\end{figure}

Figures~\ref{fig:bookday} and~\ref{fig:windowzoom} show the result. In the securities with
derivatives the spread is slightly wider through much of the expiry day and widens further in
the window, where the displayed depth also falls below that of the day before. In the placebo
group the two days are almost indistinguishable in the window. Table~\ref{tab:bookparts}
quantifies the comparison. Before the window the spread on expiry days is wider by
\BookdayDayLiquidSpread{} per cent in the liquid group and depth does not change; in the window
the spread is wider by \BookdayWindowLiquidSpread{} per cent in the liquid group and
\BookdayWindowIlliquidSpread{} per cent in the illiquid group, and depth is lower by
\BookdayWindowLiquidDepth{} and \BookdayWindowIlliquidDepth{} per cent. In the placebo group the
changes in the window are \BookdayWindowPlaceboSpread{} and \BookdayWindowPlaceboDepth{} per
cent, and neither is significant. Figure~\ref{fig:bookratio} traces the difference through the
session.

\begin{table}[htbp]
\centering
\small
\caption{Expiry day against the day before, in per cent (one hundred times the difference in
logarithms), averaged over security-months; paired $p$-values in brackets. Realised variance
is the sum of squared five-second changes of the logarithmic midpoint, per minute.}
\label{tab:bookparts}
\begin{tabular}{llccc}
\toprule
 & & Quoted spread & Depth & Realised variance \\
\midrule
Before the window & Liquid & \BookdayDayLiquidSpread{} [\BookdayDayLiquidSpreadP{}] & \BookdayDayLiquidDepth{} [\BookdayDayLiquidDepthP{}] & \BookdayDayLiquidRv{} [\BookdayDayLiquidRvP{}] \\
 & Illiquid & \BookdayDayIlliquidSpread{} [\BookdayDayIlliquidSpreadP{}] & \BookdayDayIlliquidDepth{} [\BookdayDayIlliquidDepthP{}] & \BookdayDayIlliquidRv{} [\BookdayDayIlliquidRvP{}] \\
 & Placebo & \BookdayDayPlaceboSpread{} [\BookdayDayPlaceboSpreadP{}] & \BookdayDayPlaceboDepth{} [\BookdayDayPlaceboDepthP{}] & \BookdayDayPlaceboRv{} [\BookdayDayPlaceboRvP{}] \\
\midrule
Settlement window & Liquid & \BookdayWindowLiquidSpread{} [\BookdayWindowLiquidSpreadP{}] & \BookdayWindowLiquidDepth{} [\BookdayWindowLiquidDepthP{}] & \BookdayWindowLiquidRv{} [\BookdayWindowLiquidRvP{}] \\
 & Illiquid & \BookdayWindowIlliquidSpread{} [\BookdayWindowIlliquidSpreadP{}] & \BookdayWindowIlliquidDepth{} [\BookdayWindowIlliquidDepthP{}] & \BookdayWindowIlliquidRv{} [\BookdayWindowIlliquidRvP{}] \\
 & Placebo & \BookdayWindowPlaceboSpread{} [\BookdayWindowPlaceboSpreadP{}] & \BookdayWindowPlaceboDepth{} [\BookdayWindowPlaceboDepthP{}] & \BookdayWindowPlaceboRv{} [\BookdayWindowPlaceboRvP{}] \\
\bottomrule
\end{tabular}
\end{table}

\begin{figure}[htbp]
  \centering
  \includegraphics[width=\textwidth]{bh2_ratio_by_time.pdf}
  \caption{Difference between the expiry day and the day before in the quoted spread and in
    displayed depth, in per cent, by quarter hour of the session, with 95\% intervals.}
  \label{fig:bookratio}
\end{figure}

\subsection{Why the window is costlier}

The last column of Table~\ref{tab:bookparts} points to the reason. The realised variance of the
midpoint in the window is higher on expiry days by \BookdayWindowLiquidRv{} per cent in the
liquid group and \BookdayWindowIlliquidRv{} per cent in the illiquid group, against
\BookdayWindowPlaceboRv{} per cent, not significant, in the placebo group. Liquidity suppliers
quote wider and smaller when prices move more, because their inventory is exposed to larger
price changes and they are more likely to trade with better-informed counterparties. If the
widening is this response, a security whose volatility rose more on the expiry day should show
a larger widening, and a security whose volatility did not change should show little.

Across the security-months of the derivatives groups, the change in the logarithm of the
spread in the window is regressed on the change in the logarithm of realised variance. The
slope is \VoltestFoSpreadSlope{} (standard error \VoltestFoSpreadSlopeSe{}), and the intercept,
the widening of a security whose volatility did not change, is \VoltestFoSpreadIntercept{} per
cent ($p = \VoltestFoSpreadInterceptP{}$) against an average widening of
\VoltestFoSpreadRaw{} per cent. The rise in volatility thus accounts for about
\VoltestFoSpreadExplained{} per cent of the widening. For depth the intercept is
$\VoltestFoDepthIntercept{}$ per cent ($p = \VoltestFoDepthInterceptP{}$) against an average
fall of $\VoltestFoDepthRaw{}$ per cent, so the fall in depth is accounted for by volatility
entirely. The window on expiry days is costlier to trade in because prices move more in it, and
liquidity suppliers respond to that in the usual way. The additional volatility is itself
consistent with the additional trading that the expiry draws into the window
(Section~\ref{sec:bh1}), which is two-sided rather than directional (Section~\ref{sec:bh3}).


\input{generated/tables/contrasts}
